% Part: first-order-logic
% Chapter: introduction
% Section: sentences

\documentclass[../../../include/open-logic-section]{subfiles}

\begin{document}

\olfileid{fol}{int}{snt}

\olsection{\usetoken{P}{sentence}}

صار لنا الآن رسم لتعريف الإرضاء، أي «الصدق في»، يتناول
ال!!{structure}s وال!!{formula}s. لكنه يعتمد على أمر زائد، هو
إسناد !!a{variable}، ونحن إنما أردنا تعريف ال!!{sentence}s.
فما ال!!{sentence}s، وكيف يرتفع الاحتياج إلى الإسناد؟

من مبادئ المنطق الصوري علاقة ال!!{variable}s بالمكممات. فالمكمم
يتلوه !!a{variable}، ويقيّد هذا ال!!{variable} في جزء
ال!!{sentence} الذي يتسلط عليه، وهو «نطاقه». وأما تعريف
ال!!{formula}s فلم يوجب أن يكون لكل !!{variable} مكمم يقيّده؛
فال!!{variable}s التي لا مكمم لها حرة، أو غير مقيدة. ونجعل
ال!!{sentence}s صنف ال!!{formula}s الذي يخلو من
ال!!{variable}s الحرة.

وقد يسهل إدراك متى يكون ورود !!a{variable} في
!!a{formula}~$!A$ مقيدًا، وما المكمم الذي يقيده، ومتى يكون حرًا.
وربما عرف القارئ سبيلًا لاختبار ذلك بعدّ الأقواس ونحوه. غير أن
المطلوب هنا تعريف دقيق. ولما كان تعريف ال!!{formula}s استقرائيًا،
أمكن أن نحد بالطريق نفسه ورود !!a{variable}~${\OLId{latin-ordinary}{x}}$ في
!!a{formula}~$!A$، أحر هو أم مقيد. ففي لغتنا المبسطة نضع مثلًا:
\begin{enumerate}
  \item إن كانت $!A$ ذرية، فكل ورود لـ~${\OLId{latin-ordinary}{x}}$ فيها حر؛ ومن ذلك ورود
  ${\OLId{latin-ordinary}{x}}$ في~$\Atom{\Obj P}{x}$.
  \item وإن كانت $!A$ هي $\lnot !B$، فحرية ورود~${\OLId{latin-ordinary}{x}}$ في
  $\lnot !B$ مكافئة لحرية الورود المقابل لـ~${\OLId{latin-ordinary}{x}}$ في~$!B$.
  فالورود الحر في~$!B$ يبقى حرًا في موضعه المقابل من~$\lnot !B$.
  \item وإن كانت $!A$ هي $(!B \land !C)$، فحرية ورود~${\OLId{latin-ordinary}{x}}$ في
  $(!B \land !C)$ مكافئة لحرية الورود المقابل لـ~${\OLId{latin-ordinary}{x}}$
  في~$!B$ أو في~$!C$.
  \item وإن كانت $!A$ هي $\lexists[{\OLId{latin-ordinary}{x}}][!B]$، امتنع أن يكون ورود
  لـ~${\OLId{latin-ordinary}{x}}$ في~$!A$ حرًا. أما إن كانت $\lexists[{\OLId{latin-ordinary}{y}}][!B]$، و${\OLId{latin-ordinary}{y}}$
  !!{variable} غير~${\OLId{latin-ordinary}{x}}$، فإن ورود~${\OLId{latin-ordinary}{x}}$ في $\lexists[{\OLId{latin-ordinary}{y}}][!B]$
  يكون حرًا إذا وفقط إذا كان الورود المقابل لـ~${\OLId{latin-ordinary}{x}}$ حرًا في~$!B$.
\end{enumerate}

بعد هذا الضبط للورود الحر والمقيد يتم التعريف: ال!!{sentence}
!!{formula} ليس فيها ورود حر لل!!{variable}s.

\end{document}
